% Source: https://artofproblemsolving.com/community/c3h1986756_set_of_integers
% Stable identity: aops-topic-1986756-q01
Write $(x-y)(x+y)=1988$. The factors $x-y$ and $x+y$ have the same parity (their sum is $2x$), and their product is even, so both are even. Put $x-y=2m$ and $x+y=2n$; then $mn=\dfrac{1988}{4}=497=7\cdot71$, and $x=m+n$, $y=n-m$.

The factorisations $mn=497$ are $(m,n)=(\pm1,\pm497),(\pm497,\pm1),(\pm7,\pm71),(\pm71,\pm7)$ (same signs). They give $(x,y)=(\pm498,\pm496)$ and $(\pm78,\pm64)$. So the solutions are
\[
\boxed{(x,y)=(\pm498,\pm496)\ \text{and}\ (\pm78,\pm64)}
\]
with all four sign combinations in each, $8$ pairs in total. (Check: $498^2-496^2=2\cdot994=1988$ and $78^2-64^2=14\cdot142=1988$.)
